\chapter{Requirement Analysis and Specification}

\section{Stakeholders}
The requirements in this chapter were drawn from the people and organisations that use, run or are affected by ScamGuard. Table~\ref{tab:stake} lists them with what each needs from the system.

\begin{table}[H]
\centering
\caption{Stakeholders and their needs}
\label{tab:stake}
\begin{tabular}{|L{4.0cm}|L{4.7cm}|L{5.2cm}|}
\hline
\textbf{Stakeholder} & \textbf{Role} & \textbf{What they need from ScamGuard} \\
\hline
Visitors (students, parents, older relatives, first-time UPI users) & Paste a doubtful message into the quick check without an account & A clear verdict, the trap phrases and the next safe step \\
\hline
Registered users & Run full analyses and keep a history & Analysis of text, links, e-mails, screenshots, PDFs and QR codes; history, cases, reports and analytics \\
\hline
Administrators & Run the service & Registered model versions, a robustness test and an audit trail \\
\hline
Project team & Build, test and maintain the system & Reproducible training, versioned models and automated tests \\
\hline
Guide and reviewers & Assess the research claims & Measured, repeatable results with stated limits; the cross-domain study \\
\hline
Reporting institutions (1930 helpline, cybercrime.gov.in, Sanchar Saathi) & Receive fraud reports from citizens & Users who reach them early and arrive with the evidence of the scan \\
\hline
Dataset authors & Publish the IMC 2025, UCI and India SMS collections & Use of their data within its licence, with credit \\
\hline
\end{tabular}
\end{table}

\section{Functional Requirements}\label{sec:funcreq}
Functional requirements describe what the system must be able to do.
\begin{enumerate}
\item \textbf{Quick check.} Without creating an account, a visitor pastes a message on the home page and gets a verdict, a risk score, the trap phrases highlighted and a short explanation.
\item \textbf{Accounts.} People can register and sign in (with a captcha) or try a demo account, and they can edit, export or delete their own data.
\item \textbf{Full analysis.} Once signed in, a user can submit text, a URL, an e-mail, a screenshot, a PDF or a QR code. The result shows the verdict, the classifier's probability, the final risk score, the Scam X-Ray stages, the TrustLens report and suggested actions.
\item \textbf{Explanation.} Every result lists the phrases that matched, the scam stage each belongs to (hook, pressure, request, cash-out), the evidence on both sides and the way a genuine message would look.
\item \textbf{Cross-domain page.} A visitor picks the baseline model (UCI only) or the adapted one (UCI plus 40\% of the Indian pool), chooses Naive Bayes or Logistic Regression and classifies a message. The page returns the label, the spam probability, the model version and the features that drove the decision, next to the study's measured results.
\item \textbf{AI copilot.} Follow-up questions go to a hosted language model. Each visitor has a limit, there is a daily cap, and fixed fallback answers cover outages.
\item \textbf{Safety tools.} A QR code check, a ``before you pay'' check for UPI requests, a check for payment screenshots, a file check that catches disguised Android apps, and a chat timeline.
\item \textbf{History, cases and reports.} Past scans can be searched and grouped into cases; users can leave feedback, download a PDF report and post a scam to the community view.
\item \textbf{Analytics and administration.} Users see their totals and trends. Administrators also see the registered model versions and can run a robustness test.
\end{enumerate}

\section{Non-Functional requirements}
These requirements describe the qualities the system must have.
\begin{itemize}
\item \textbf{Performance.} On a two-core server the classifier needs about 2~ms per message, about 4~ms with the intent layer. A complete web request takes roughly 1--2 seconds.
\item \textbf{Reliability.} Models are read from versioned files once, when the server starts. The back end keeps itself from sleeping on the free tier, and the front end retries short-lived proxy errors.
\item \textbf{Security.} Passwords are stored as bcrypt hashes; access relies on signed JWT tokens whose refresh tokens rotate; every endpoint has a rate limit and a size limit on the request body; fetching outside URLs is protected against SSRF; and security headers are set.
\item \textbf{Privacy.} Nothing typed into the quick check or the cross-domain page is saved, and contributed messages are masked before they are stored.
\item \textbf{Usability.} The interface has an elder mode, can be driven from the keyboard and fits phone screens.
\item \textbf{Maintainability and reproducibility.} Training and serving call the same preprocessing code. Continuous integration runs 836 back-end and ML tests and 14 front-end tests, and a single notebook plus a single script regenerate every number we report.
\end{itemize}

\section{System Requirements}
ScamGuard needs no GPU. Table~\ref{tab:hw} lists the hardware each part of the system needs and Table~\ref{tab:sw} the software used for each layer.

\begin{table}[H]
\centering
\caption{Hardware requirements}
\label{tab:hw}
\begin{tabular}{|L{4.0cm}|L{10.3cm}|}
\hline
\textbf{Component} & \textbf{Requirement} \\
\hline
Development machine & An ordinary laptop with 8~GB RAM and a recent 64-bit CPU; a GPU is not required \\
\hline
Training & On a laptop CPU, under 3 seconds for the production model and around a minute for the whole cross-domain notebook \\
\hline
Server & A single small cloud instance; 512~MB RAM is enough \\
\hline
Client & Any web browser on a phone or a computer \\
\hline
\end{tabular}
\end{table}

\begin{table}[H]
\centering
\caption{Software requirements}
\label{tab:sw}
\begin{tabular}{|L{4.0cm}|L{10.3cm}|}
\hline
\textbf{Layer} & \textbf{Software} \\
\hline
Language & Python 3.10+ (tested on 3.13), TypeScript 5.5 \\
\hline
Machine learning & scikit-learn 1.7.2, NumPy, pandas, SciPy, joblib, datasketch, ftfy, matplotlib \\
\hline
Notebook and study app & Jupyter or Google Colab, Streamlit \\
\hline
Back end & FastAPI 0.115, Uvicorn, SQLAlchemy 2.0, Alembic, Pydantic 2, SlowAPI, python-jose, passlib, ReportLab, Tesseract OCR \\
\hline
Front end & Next.js 14.2, React 18.3, Tailwind CSS 3.4, Recharts, Vitest \\
\hline
Database & PostgreSQL in production, SQLite locally \\
\hline
Deployment & Free cloud tiers (front end and back end), GitHub Actions (continuous integration) \\
\hline
\end{tabular}
\end{table}

\textbf{Cost estimation.} ScamGuard was built as a student project on free tiers and open-source software, so the cost of building and running it at the present scale is zero. Python, scikit-learn, FastAPI, Next.js, Streamlit, Jupyter and VS Code are open source and cost nothing. All three public datasets (the IMC 2025 smishing reports, the UCI SMS Spam Collection and the India Spam SMS collection) were used under their open licences and cost nothing. The back end and the database run on free cloud plans. The copilot calls a hosted language model within its free quota; a daily cap and a per-visitor limit keep it inside that quota. The total estimated cost is zero for development and hosting.

\section{Data Requirements}
The work needs three kinds of labelled text, each with a clear licence:
\begin{itemize}
\item \textbf{Indian scam and genuine SMS for the main detector.} Both classes must be present and real: reported scams alone teach a model nothing about genuine bank alerts. Report labels must be checked before use, and near-identical messages must be grouped so that a campaign never sits on both sides of the split. This is met by ScamGuard-IN-2026 (11,399 messages; Section~\ref{sec:data}).
\item \textbf{A source and a target domain for the cross-domain study.} The source is the UCI SMS Spam Collection (5,574 messages, CC BY 4.0) and the target is the India Spam SMS Classification set (2,267 messages, MIT licence), both named by our guide.
\item \textbf{An untouched final test set} in each study. It is never used for training, for fitting the TF-IDF vocabulary or for choosing settings.
\end{itemize}
Personal details in stored and contributed messages (phone numbers, account numbers, e-mail addresses) must be masked. Synthetic messages may be used for training only and must be marked as synthetic in the data. The fields of each dataset are given in Section~\ref{sec:datadesign}.

\section{Constraints and Limitations}
\begin{itemize}
\item \textbf{Hosting.} The free-tier hosting allows 512~MB of memory, so large transformer models cannot be served in production; the detector must be a small model that runs on a CPU.
\item \textbf{Time span of the data.} Most reported messages in ScamGuard-IN-2026 date from 2022--2024. Scam types from 2025--26 enter training only through synthetic messages and enter testing only through 12 advisory messages.
\item \textbf{Annotation.} A single annotator made the manual audit decisions, so agreement between annotators cannot be reported.
\item \textbf{Input.} The system judges single messages only; it does not see the sender or earlier messages in the thread.
\item \textbf{Domains.} The cross-domain study has one source and one target dataset, and three seeds give only a rough idea of the spread.
\end{itemize}

\section{Assumptions}
\begin{itemize}
\item The audited labels are taken to be right, and where we could not tell we dropped the message instead of guessing.
\item Each public dataset is taken to be a fair enough sample of its domain for a controlled comparison.
\item The final target test set is kept out of training, feature fitting and every choice of settings.
\item Users paste the message as they received it, and the verdict is read as guidance, not as proof that a crime took place.
\end{itemize}

\section{Risks and Challenges}
The main risks and how they are handled are listed in Table~\ref{tab:risks}. The hardest challenge was the first one: a model trained on reports alone looked accurate but flagged one genuine message in three.

\begin{table}[H]
\centering
\caption{Risks and how they are handled}
\label{tab:risks}
\begin{tabular}{|L{4.6cm}|L{9.7cm}|}
\hline
\textbf{Risk} & \textbf{Mitigation} \\
\hline
Reported data carries wrong labels & Audit by rule and by hand, then a second pass on cases a cross-validated model disputes \\
\hline
Near-duplicates make scores look better than they are & Split by campaign; group-aware split check in the cross-domain study \\
\hline
New scam types are missing from training data & Synthetic training messages, intent rules, and a small labelled local sample (adaptation) \\
\hline
Disguised spellings (K.Y.C, bl0cked) & Caught by the intent rules; normalising them before features are built is left for future work \\
\hline
The free-tier server falls asleep & Keep-alive pings, automatic retries \\
\hline
The outside AI service is slow or down & Time limits and fallback answers; the detector itself never depends on it \\
\hline
\end{tabular}
\end{table}
